\section{Introduction}

Sparse MoE models activate only a small subset of feed-forward experts for each
token. This reduces arithmetic relative to dense parameter count, but makes
serving depend on a dynamic expert working set. Systems that offload, replicate,
or redistribute experts would benefit from knowing this set before a token
reaches the router.

Agent workloads provide an unusual source of early information. Before issuing
an LLM request, an orchestrator may know the node role, trajectory phase, prior
tool outcome, and prompt-block structure. These fields are available before
model execution and can therefore influence systems decisions without reading
the current request's hidden states or router output.

\project asks whether this context predicts expert demand beyond global
popularity and recent routing history. Its design follows a strict separation:
the predictor emits advisory working sets, while the native model router alone
determines execution and output.

This paper develops three foundations for answering the question:
\begin{itemize}[leftmargin=*]
  \item a request contract that represents the state immediately before target
        generation and excludes target answers or current agent actions;
  \item a real vLLM capture path for aligned expert IDs and router scores; and
  \item an online, source-group-isolated evaluation of global, temporal, and
        metadata-conditioned predictors.
\end{itemize}

We intentionally defer prefetch, scheduling, and placement integration until
the metadata signal passes this evidence gate.
